\section{课程总结与前沿展望}

\subsection{课程内容回顾}

KAIST CS492D 在本学期重点覆盖了扩散模型与 Flow 模型的理论与应用，课程内容可以分为两大部分：

\subsubsection{理论基础}

\begin{enumerate}
    \item \textbf{DDPM}（Denoising Diffusion Probabilistic Models）：扩散模型的基础框架，定义了前向加噪过程和反向去噪过程
    \item \textbf{DDIM}（Denoising Diffusion Implicit Models）：确定性采样方法，加速了推理过程
    \item \textbf{连续时间扩散}：将离散的扩散过程推广到连续时间 SDE/ODE 框架
    \item \textbf{SDE 与 ODE 的转换}：建立了随机与确定性采样之间的联系
    \item \textbf{Flow Matching}：一种更灵活的生成建模方法，将扩散模型连接到最优传输理论
\end{enumerate}

\subsubsection{应用拓展}

\begin{enumerate}
    \item \textbf{推理时引导}（Inference-Time Guidance）：在不重新训练模型的情况下，通过引导信号控制生成过程
    \item \textbf{Score Distillation}：将预训练扩散模型的知识蒸馏到其他数据域，实现跨模态生成
    \item \textbf{3D 生成}：利用 2D 图像扩散模型的先验知识生成 3D 内容
\end{enumerate}

\begin{importantbox}{课程设计理念}
本学期的课程设计注重从基础到应用的渐进式推进：先理解模型的数学本质（概率论、随机微分方程、最优传输），再探讨如何将这些理论工具应用到实际问题中（图像编辑、3D 生成、跨模态迁移）。
\end{importantbox}

\subsection{未来方向：离散扩散与多模态模型}

课程最后，Minhyuk Sung 教授提出了一个发人深省的问题：

\begin{importantbox}{Autoregressive vs. Diffusion：边界正在模糊}
传统观点认为：
\begin{itemize}
    \item \textbf{自回归模型}（Autoregressive Models）最适合文本等\textbf{离散序列数据}
    \item \textbf{扩散/Flow 模型}最适合图像、视频等\textbf{高维连续数据}
\end{itemize}

但这种分界正在被打破：
\begin{itemize}
    \item 自回归模型通过图像 tokenization（如 VQ-VAE）也开始生成高质量图像
    \item 离散扩散模型（Discrete Diffusion）将扩散过程推广到离散数据空间
    \item GPT-4o、Gemini 等多模态模型可以同时生成文本和图像
\end{itemize}
\end{importantbox}

\subsubsection{离散扩散模型}

扩散模型不局限于连续数据，也可以应用于离散数据。离散扩散的前向过程不是添加高斯噪声，而是通过\textbf{离散马尔可夫链}逐步破坏数据（如随机替换 token）。这为将扩散模型应用于文本、代码、分子序列等打开了大门。

\subsubsection{统一的多模态生成}

当前的前沿研究正在探索将文本和图像统一到同一个生成框架中：

\begin{knowledgebox}{多模态统一架构的两种范式}
\begin{enumerate}
    \item \textbf{Tokenize Everything}：将图像离散化为 token（通过 VQ-VAE 等），然后与文本 token 一起用自回归模型处理。代表：GPT-4o 系列。
    \item \textbf{Hybrid Architecture}：文本部分使用自回归解码，图像部分使用扩散/Flow 解码，通过共享的表示空间连接。代表：Transfusion。
\end{enumerate}

哪种范式会胜出仍是开放问题。
\end{knowledgebox}

\subsection{对研究者的建议}

Minhyuk Sung 教授在课程结束时给出了几条建议：

\begin{enumerate}
    \item \textbf{动手实践}：SDS 等技术已有大量开源实现，建议亲自运行代码、观察结果、理解失败模式
    \item \textbf{不仅关注生成，也关注编辑}：3D 编辑（在已有 3D 内容上进行修改）可能比从零生成更具实用价值
    \item \textbf{跨领域思考}：SDS 的框架不限于视觉内容，音频、分子、运动等领域都有潜在应用
    \item \textbf{关注前沿趋势}：自回归与扩散的融合、多模态统一模型等方向值得持续关注
\end{enumerate}

\subsection{把本讲内容转成研究计划}

如果把这节课当成一个研究起点，而不是单纯的综述，最关键的是把 ``SDS 能做什么'' 进一步拆成可执行的实验路线。很多文本到 3D 项目失败，不是因为没有大模型先验，而是没有把\textbf{表示、优化、评测}三件事同时设计清楚。

\begin{knowledgebox}{SDS 项目的最小实验设计}
\begin{table}[H]
\centering
\begin{tabular}{p{0.16\textwidth} p{0.28\textwidth} p{0.34\textwidth}}
\toprule
模块 & 最小可行选择 & 需要优先记录的观测 \\
\midrule
3D 表示 & NeRF 或 3D Gaussian Splatting & 视角一致性、几何平滑度、训练稳定性 \\
2D 先验 & Stable Diffusion / Imagen 类文本条件扩散模型 & prompt 敏感性、颜色饱和度、细节恢复情况 \\
优化策略 & SDS baseline + regularization + camera sampling & Janus 频率、面片漂移、收敛速度 \\
评测方式 & 多视图渲染检查 + CLIP/用户打分 + 几何可打印性 & 结果是否只是 ``看起来像''，还是具备真实 3D 一致性 \\
\bottomrule
\end{tabular}
\caption{从课程内容出发搭建首个文本到 3D 研究原型}
\end{table}
\end{knowledgebox}

\begin{warningbox}{常见误判：2D 质量高不等于 3D 质量高}
单张渲染图很漂亮，并不能说明模型真的学到了稳定几何。SDS 体系里最容易出现的误判，是研究者被某几个视角的高保真图像说服，却忽略了绕到背面后结构坍塌、纹理重复、物体拓扑错误等问题。因此实验记录必须包含连续视角巡航，而不是只展示最优样例。
\end{warningbox}

\subsection{失败模式与调参抓手}

从研究实践看，SDS 项目最难的部分不是 ``把 loss 跑起来''，而是识别当前失败究竟来自哪一层：文本条件不够清晰、2D 扩散先验过强、camera sampling 不合理，还是 3D 表示本身的容量不足。把这些问题混在一起，会导致实验只剩下盲目试 prompt。

\begin{table}[H]
\centering
\begin{tabular}{p{0.22\textwidth} p{0.3\textwidth} p{0.28\textwidth}}
\toprule
现象 & 更可能的根因 & 优先尝试的修复动作 \\
\midrule
正面像样、背面崩塌 & 多视角约束不足，采样视角分布过窄 & 扩大 elevation/azimuth 覆盖范围，引入 view-dependent regularization \\
颜色过饱和、纹理黏连 & SDS 梯度过强，2D 先验牵引过度 & 调低 guidance 强度，增加 entropy/TV 正则 \\
几何浮肿、边界模糊 & 表示容量不足或 density regularization 不稳 & 更换 3D 表示、增加几何先验、显式约束 opacity \\
结果高度依赖 prompt 细节 & 文本条件承担了过多结构信息 & 改写 prompt 模板，并配合参考视角或 weak supervision \\
\bottomrule
\end{tabular}
\caption{SDS 文本到 3D 项目中常见失败模式与修复动作}
\end{table}

\begin{importantbox}{课程收官留下的真正方法论}
这门课最后一讲最重要的收获不是 ``又学会一个新算法''，而是理解扩散模型研究正在从单纯的图像生成，转向\textbf{把强先验迁移到稀缺数据域}。SDS、VSD、离散扩散、多模态统一架构都属于这个更大的母问题：如何让一个在超大数据上学到的生成器，为另一个数据稀缺但价值很高的表示空间提供优化信号。
\end{importantbox}

\subsection{本章小结}

课程系统地覆盖了从 DDPM 到 Flow Matching 的理论发展，以及从推理引导到 Score Distillation 的应用拓展。当前的前沿趋势——离散扩散、自回归图像生成、多模态统一模型——正在打破传统的模型选择边界，预示着更灵活、更强大的生成建模框架。

\newpage
